\section*{Chapitre \ref{ch:pl:pricing-library-architecture} --- Architecture de la bibliothèque de pricing}

\begin{solution}{exo:pl:pricing-library-architecture:1}
\texttt{at expiry \{ if S > 105 \{ pay 10; \} \}} avec le calendrier \texttt{expiry} contenant la date d'échéance. Sur le marché du chapitre, son prix s'établit à
4.09, égal à celui de l'option digitale de la bibliothèque valorisée par son moteur Monte Carlo sur les mêmes trajectoires.
\end{solution}

\begin{solution}{exo:pl:pricing-library-architecture:2}
L'autocallable est comparé au moteur Monte Carlo de la bibliothèque sur les mêmes trajectoires (même identifiant, même graine, mêmes dates) et à la même
arithmétique de payoff~: les nombres sont identiques. L'européenne est comparée à la formule analytique, qui n'a aucune erreur Monte Carlo~; l'estimation du script
diffère d'elle de sa propre erreur d'échantillonnage, soit 0.82 erreur type ici.
\end{solution}

\begin{solution}{exo:pl:pricing-library-architecture:3}
Parce que la fiche descriptive verse le coupon à une date d'observation puis rembourse si le titre est rappelé~: si l'autocall passe en premier, \texttt{stop}
mettrait fin à la trajectoire avant le paiement du coupon de la date de rappel, et le script sous-évaluerait chaque trajectoire rappelée d'un coupon (et de sa mémoire).
\end{solution}

\begin{solution}{exo:pl:pricing-library-architecture:4}
Elle retire les trajectoires couvertes par le masque courant de l'ensemble vivant~; chaque instruction et chaque événement ultérieur s'exécute sous \texttt{mask \& alive}, de sorte que ces
trajectoires ne paient plus rien et ne changent plus rien. Sauter la suite du bloc reviendrait à le sauter pour toutes les trajectoires, y compris celles auxquelles le \texttt{stop}
ne s'appliquait pas.
\end{solution}

\begin{solution}{exo:pl:pricing-library-architecture:5}
L'évaluateur compilé paie un coût fixe par événement (parcours de l'arbre une seule fois) et un petit coût par trajectoire~; l'interpréteur paie le parcours à chaque trajectoire. Pour
peu de trajectoires, le coût fixe et la génération des trajectoires (commune aux deux) pèsent davantage~; à mesure que le nombre de trajectoires croît, le rapport tend vers le rapport des coûts par trajectoire, et
se stabilise dès que les deux sont dominés par le travail sur les tableaux et le trafic mémoire.
\end{solution}

\begin{solution}{exo:pl:pricing-library-architecture:6}
Des observateurs pour l'écran, où une seule cotation de marché change à la fois et où seuls les prix affichés doivent être recalculés~; des instantanés pour le run de risque,
où des milliers de marchés bumpés sont évalués indépendamment et en parallèle, et où rien ne doit dépendre de l'ordre des mises à jour.
\end{solution}

\begin{solution}{exo:pl:pricing-library-architecture:7}
97.64 à 50~\%, 94.25 à 60~\%, 92.28 à 70~\%~: environ 5.4 points sur vingt points de barrière. Un knock-in lookback est très sensible à la
barrière, puisque toute observation située sous elle compte~; le niveau de barrière est un paramètre de valorisation à négocier et un risque à couvrir.
\end{solution}

\begin{solution}{exo:pl:pricing-library-architecture:8}
Des nombres aléatoires indépendants font différer le script et la bibliothèque par le bruit Monte Carlo, ce qui masque de petites erreurs de payoff~; ils renoncent
aux nombres aléatoires communs, de sorte que les prix bumpés du script sont bruités et les grecques instables. Il faut tirer les trajectoires de la bibliothèque avec la graine de l'instrument~; recouper
par comparaison avec les moteurs sur les mêmes trajectoires (exact) et avec les formules fermées (à l'erreur type près).
\end{solution}

\begin{solution}{pb:pl:pricing-library-architecture:1}
\begin{enumerate}
\item Instruments, \omterm{def:pl:pricing-library-architecture:market}{objets de marché}, modèles et moteurs~; la fonction de repricing sur les instantanés les relie.
\item Questions portant sur un élément de l'état du marché (un facteur d'actualisation, une volatilité)~; un instantané ne change jamais, de sorte qu'un bump en crée un nouveau et
rien n'est périmé.
\item Nombres aléatoires communs~: la graine ne dépend que de l'instrument, de sorte que les prix de base et les prix bumpés utilisent les mêmes trajectoires.
\item \texttt{supports} et \texttt{price} pour un couple (instrument, modèle).
\item Parce que chaque grecque, chaque scénario et chaque lot est cette fonction appliquée à des instantanés bumpés.
\item État, événements sur calendrier, assignations, conditions, paiements et arrêts~; \texttt{S}, \texttt{perf}, \texttt{t}, les paramètres et
l'état.
\item À partir des blocs \texttt{at}~: chaque date de chaque bloc, triée par date, dans l'ordre du script à l'intérieur d'une même date.
\item L'interpréteur parcourt l'arbre trajectoire par trajectoire avec des scalaires~; l'évaluateur compilé le parcourt événement par événement avec des tableaux et des masques.
\item Une variable d'état pour la performance la plus basse, mise à jour à chaque observation, et le test d'échéance qui s'appuie sur elle.
\item L'analyseur signale une erreur avec sa ligne et sa colonne.
\item Européenne 11.3953 contre 11.3485 (0.82)~; asiatique 6.8012 contre 6.8012 (0.00)~; down-and-out 11.0373 contre 11.0226 (0.26)~;
autocallable 95.8750 contre 95.8750 (0.00)~; lookback 94.2474, sans classe correspondante dans la bibliothèque.
\item Les deux sont comparés au moteur Monte Carlo de la bibliothèque sur les mêmes trajectoires et payoffs.
\item 0.1518 par unité de spot pour 100 de notionnel, pour les deux.
\item 18.8 fois à 1\,000 trajectoires, 33.4 à 30\,000, 34.3 au plus (à 10\,000).
\item Deux calculs pour quatre états du marché~: un à la première lecture, un à la lecture après trois changements.
\item \textbf{Résultat nommé.} L'autocallable scripté se valorise à 95.8750 pour 100 avec un delta de 0.1518, exactement comme le moteur autocallable de la bibliothèque
sur nombres aléatoires communs (à 0.00 erreur type près)~; le script compilé s'exécute 22 à 37 fois plus vite que le script interprété.
\item Validation par rapport aux moteurs de la bibliothèque et aux formules fermées, revue du script, son stockage avec une version, et approbation du modèle.
\item Lorsqu'il a besoin d'une méthode numérique que le moteur de script n'offre pas (exercice anticipé par régression, une EDP), ou lorsqu'il se traite en volume
et mérite un moteur dédié, plus rapide.
\item Sous contrôle de version en texte, la transaction renvoyant à une version du script et à ses paramètres~; chaque modification est revue et repricée.
\item De nouveaux produits décrits comme des données, valorisés et gérés en risque par le dispositif auquel la bibliothèque fait déjà confiance.
\end{enumerate}
\end{solution}

\begin{solution}{iq:pl:pricing-library-architecture:1}
Pour que chacun varie indépendamment~: un nouveau modèle valorise les anciens produits, un nouveau moteur accélère les anciens modèles, un nouveau produit utilise les moteurs existants~;
le risque et les scénarios fonctionnent sur n'importe quelle combinaison grâce à une seule fonction de repricing.
\iqlookfor{Séparation des responsabilités, un seul chemin de repricing.}
\end{solution}

\begin{solution}{iq:pl:pricing-library-architecture:2}
On l'écrit comme un \omterm{def:pl:pricing-library-architecture:script}{script de payoff} sur le moteur de scripts de la bibliothèque, on le valide face aux produits les plus proches que la bibliothèque valorise, on vérifie les grecques et
les limites, on le fait relire~; un nouveau moteur pourra venir plus tard si le volume le justifie.
\iqlookfor{Scripting et validation.}
\end{solution}

\begin{solution}{iq:pl:pricing-library-architecture:3}
Évaluer sur toutes les trajectoires à la fois avec des tableaux~: les conditions deviennent des masques, les paiements des sommes masquées, la terminaison un ensemble vivant~; générer les trajectoires une seule fois~;
éviter les boucles Python par trajectoire, et revenir au code compilé si le parcours de l'arbre devient lui-même le coût.
\iqlookfor{Vectorisation par masques.}
\end{solution}

\begin{solution}{iq:pl:pricing-library-architecture:4}
Résultats marqués périmés lors des changements d'entrée et recalculés à la lecture~; mauvais lorsque les lectures sont constantes (aucune économie), dans un risque par lots en parallèle (état
mutable partagé), et lorsque les chaînes de notification s'allongent et que leur ordre compte.
\iqlookfor{Observateurs face aux instantanés immuables.}
\end{solution}

\begin{solution}{iq:pl:pricing-library-architecture:5}
Une petite grammaire (événements sur calendrier, état, conditions, paiements, arrêts), un analyseur avec positions d'erreur, l'extraction du calendrier, un évaluateur sur
les trajectoires de la bibliothèque avec nombres aléatoires communs, l'enregistrement comme moteur, une suite de validation face aux produits de la bibliothèque, un stockage versionné et
la revue des scripts.
\iqlookfor{Petit langage, trajectoires de la bibliothèque, validation, versionnement.}
\end{solution}
